\chapter{Configurații Kubernetes}
\label{app:kubernetes}

Această anexă include manifestele declarative pentru rularea sistemului
pe un cluster Kubernetes. Toate resursele sunt aplicate în namespace-ul
\texttt{epimodel} și urmează structura arhitecturală descrisă în
capitolul~\ref{ch:arhitectura}: API stateless, frontend servit prin Nginx,
runner-e scalabile independent, Redis pentru cozi de mesaje și YugabyteDB
pentru persistență.

Manifestele se află în directorul \texttt{code/k8s/} al repository-ului și
sunt organizate cu prefixe numerice pentru aplicare secvențială
(\texttt{kubectl apply -f code/k8s/}). \Cref{tab:k8s_manifeste} sintetizează
conținutul directorului.

\begin{table}[htbp]
\centering
\caption{Manifestele Kubernetes incluse în repository.}
\label{tab:k8s_manifeste}
\renewcommand{\arraystretch}{1.2}
\begin{tabular}{p{4cm}p{4cm}p{5cm}}
\toprule
\textbf{Fișier} & \textbf{Resurse principale} & \textbf{Rol} \\
\midrule
\texttt{00-namespace.yaml}         & \texttt{Namespace}                 & Izolează aplicația într-un namespace dedicat \\
\texttt{01-configmap.yaml}         & \texttt{ConfigMap}                 & Variabile non-sensibile pentru API \\
\texttt{02-secret.example.yaml}    & \texttt{Secret} (template)         & Șablon pentru JWT și DSN-ul bazei de date \\
\texttt{10-yugabyte-master.yaml}   & \texttt{StatefulSet} ($\times 3$)  & YugabyteDB master (coordonatorul cluster-ului) \\
\texttt{11-yugabyte-tserver.yaml}  & \texttt{StatefulSet} ($\times 3$)  & YugabyteDB tserver (YSQL pe portul 5433) \\
\texttt{20-redis.yaml}             & \texttt{StatefulSet} ($\times 1$)  & Coadă Redis cu persistență AOF \\
\texttt{30-migrate-job.yaml}       & \texttt{Job}                       & Aplică migrațiile SQL prin \texttt{golang-migrate} \\
\texttt{40-api.yaml}               & \texttt{Deployment} + \texttt{Service} & Backend Go (REST API) \\
\texttt{50-runner-python.yaml}     & \texttt{Deployment}                & Runner Python (\texttt{scipy.solve\_ivp}) \\
\texttt{51-runner-r.yaml}          & \texttt{Deployment}                & Runner R (\texttt{deSolve::ode}) \\
\texttt{52-runner-julia.yaml}      & \texttt{Deployment}                & Runner Julia (\texttt{DifferentialEquations.jl}) \\
\texttt{60-frontend.yaml}          & \texttt{Deployment} + \texttt{Service} & Frontend Vue.js servit prin Nginx \\
\texttt{70-ingress.yaml}           & \texttt{Ingress}                   & Rutare HTTP externă (\texttt{/api} și \texttt{/}) \\
\bottomrule
\end{tabular}
\end{table}

Listările care urmează reproduc manifestele cele mai relevante pentru
explicitarea arhitecturii: namespace-ul de izolare, stratul de persistență
(YugabyteDB tserver), coada de mesaje (Redis), API-ul, un exemplu de runner
(Julia, cel mai mare consumator de memorie) și ingress-ul de rutare. Restul
manifestelor (master-ul YugabyteDB, ConfigMap-ul, template-ul Secret, Job-ul
de migrare, celelalte două runnere, frontend-ul) urmează aceleași convenții
și pot fi consultate direct în repository.

\section*{Namespace dedicat}

\lstinputlisting[language=bash, caption={\texttt{00-namespace.yaml} --- Namespace-ul dedicat aplicației.}]{../code/k8s/00-namespace.yaml}

\section*{YugabyteDB tserver}

Tserver-ul rulează ca \texttt{StatefulSet} cu trei replici și expune două
servicii: unul \emph{headless} pentru descoperirea internă a peer-urilor
(necesar protocolului Raft) și unul \texttt{ClusterIP} pentru consumatorii
de YSQL (API-ul și runner-ele).

\lstinputlisting[language=bash, caption={\texttt{11-yugabyte-tserver.yaml} --- StatefulSet și servicii pentru YugabyteDB tserver.}]{../code/k8s/11-yugabyte-tserver.yaml}

\section*{Redis}

Redis-ul este implementat ca \texttt{StatefulSet} cu o singură replică și
volum persistent pentru AOF (\emph{append-only file}), pentru ca starea
cozilor să supraviețuiască restartărilor pod-ului.

\lstinputlisting[language=bash, caption={\texttt{20-redis.yaml} --- Redis cu persistență AOF.}]{../code/k8s/20-redis.yaml}

\section*{API Go}

API-ul rulează ca \texttt{Deployment} cu două replici pentru toleranță la
pierderea unui pod. Configurația non-sensibilă vine din \texttt{ConfigMap},
iar parolele și secretele din \texttt{Secret}. Probele HTTP pe
\texttt{/api/health} verifică atât pod-ul, cât și conectivitatea efectivă
către baza de date și Redis.

\lstinputlisting[language=bash, caption={\texttt{40-api.yaml} --- Deployment și Service pentru API-ul Go.}]{../code/k8s/40-api.yaml}

\section*{Runner Julia}

Runner-ul Julia este Deployment-ul cu cele mai generoase limite de resurse
($512$~MiB request, $2$~GiB limit) din cauza costului de compilare JIT al
\texttt{DifferentialEquations.jl}. Runner-ele R și Python au structuri
identice, dar cu limite mai mici.

\lstinputlisting[language=bash, caption={\texttt{52-runner-julia.yaml} --- Deployment pentru runner-ul Julia.}]{../code/k8s/52-runner-julia.yaml}

\section*{Ingress de rutare}

Ingress-ul nginx ruteză cererile externe către cele două servicii vizibile
public: API-ul Go pentru cererile sub prefixul \texttt{/api} și frontend-ul
Vue.js pentru restul.

\lstinputlisting[language=bash, caption={\texttt{70-ingress.yaml} --- Ingress nginx pentru rutare HTTP.}]{../code/k8s/70-ingress.yaml}

\section*{Alternativă pentru dezvoltare locală}

În paralel cu manifestele Kubernetes, repository-ul include și o configurație
Docker Compose care reproduce aceeași topologie pe o singură mașină, fără
nevoia unui cluster. Aceasta este utilă pentru dezvoltare locală și pentru
verificări rapide ale fluxurilor end-to-end fără costul operațional al unui
cluster K8s real.

\lstinputlisting[language=bash, caption={\texttt{docker-compose.yml} --- Configurația echivalentă pentru mediul de dezvoltare local.}]{../code/docker-compose.yml}
